Attention and modern Hopfield retrieval share a familiar computation: a query
is mapped to a softmax-weighted average of stored vectors
\cite{vaswaniAttentionAllYou2017,ramsauerHopfieldNetworksAll2021}. Adding
calibrated Gaussian noise to the corresponding gradient update gives an
unadjusted Langevin algorithm (ULA), which we previously described as
\emph{stochastic attention}. That construction is mathematically valid as a
numerical dynamics, but it leaves a more basic question unanswered: what is the
probability distribution being sampled, and does it require a Markov chain?

For the standard tied-key/value modern Hopfield energy, the answer is simple.
Completing the square shows that its Boltzmann law is a finite isotropic
Gaussian mixture centered on the stored memories. Softmax attention gives the
posterior responsibilities of those components, and the attention-minus-state
vector is exactly the mixture score. An independent equilibrium endpoint is
therefore available by ancestral sampling: choose a memory according to its
analytic component weight and add Gaussian noise. When all memories are
$\ell_2$-normalized, those weights are uniform. This exact representation means
that ULA adds no distributional expressiveness to the unconstrained endpoint
problem and that a finite-step, warm-started ULA output should not be conflated
with an exact draw from the stationary law.

This correction changes the empirical question. The relevant object is a
memory-centered density closely related to an isotropic Parzen estimator, not a
general-purpose generator. Temperature sets the perturbation scale and controls
how concentrated posterior responsibilities are around individual memories.
Hard masking removes mixture components, while multiplicity biases reweight
them; neither operation by itself requires Langevin sampling. Finite-time ULA
or MALA can still be useful when trajectories, warm starts, or later constrained
extensions are of interest, but exact ancestral draws are the reference for
independent equilibrium endpoints.

This revision makes four contributions. First, we derive the normalized
Gaussian-mixture law for arbitrary memory norms and multiplicities. Second, we
identify attention weights with component responsibilities and derive the exact
score. Third, we compare exact draws with ULA and MALA on a tractable target and
reanalyze the submitted MNIST and protein experiments using all memory
components. Fourth, we delimit the empirical scope: the reported novelty,
diversity, energy, and protein statistics describe local memory-supported
behavior; they do not by themselves demonstrate realism, functional validity,
or superiority to learned generative models. This distinction between the
probability model and algorithms used to explore it is the central result.
